\documentclass[11pt]{article}

% Change "review" to "final" to generate the final (sometimes called camera-ready) version.
% Change to "preprint" to generate a non-anonymous version with page numbers.
\usepackage[review]{acl}

% Standard package includes
\usepackage{times}
\usepackage{latexsym}

% For proper rendering and hyphenation of words containing Latin characters (including in bib files)
\usepackage[T1]{fontenc}
% For Vietnamese characters
% \usepackage[T5]{fontenc}
% See https://www.latex-project.org/help/documentation/encguide.pdf for other character sets

% This assumes your files are encoded as UTF8
\usepackage[utf8]{inputenc}

% This is not strictly necessary, and may be commented out,
% but it will improve the layout of the manuscript,
% and will typically save some space.
\usepackage{microtype}

% This is also not strictly necessary, and may be commented out.
% However, it will improve the aesthetics of text in
% the typewriter font.
\usepackage{inconsolata}

%Including images in your LaTeX document requires adding
%additional package(s)
\usepackage{graphicx}

\usepackage{booktabs}
\usepackage{longtable}
\usepackage{array}
\usepackage{xcolor}
\usepackage{colortbl}
\usepackage{caption}
\usepackage{adjustbox}   % for \begin{adjustbox}
\usepackage{natbib}      % for \citeyear  (or use biblatex)
\usepackage{lscape}      % optional: landscape

% If the title and author information does not fit in the area allocated, uncomment the following
%
%\setlength\titlebox{<dim>}
%
% and set <dim> to something 5cm or larger.

\title{A Survey on the Evolving Frontier of Multi-Page Visually Rich Document Understanding: Architectures, Trends, Methods, and Challenges}

% Author information can be set in various styles:
% For several authors from the same institution:
% \author{Author 1 \and ... \and Author n \\
%         Address line \\ ... \\ Address line}
% if the names do not fit well on one line use
%         Author 1 \\ {\bf Author 2} \\ ... \\ {\bf Author n} \\
% For authors from different institutions:
% \author{Author 1 \\ Address line \\  ... \\ Address line
%         \And  ... \And
%         Author n \\ Address line \\ ... \\ Address line}
% To start a separate ``row'' of authors use \AND, as in
% \author{Author 1 \\ Address line \\  ... \\ Address line
%         \AND
%         Author 2 \\ Address line \\ ... \\ Address line \And
%         Author 3 \\ Address line \\ ... \\ Address line}

\author{First Author \\
  Affiliation / Address line 1 \\
  Affiliation / Address line 2 \\
  Affiliation / Address line 3 \\
  \texttt{email@domain} \\\And
  Second Author \\
  Affiliation / Address line 1 \\
  Affiliation / Address line 2 \\
  Affiliation / Address line 3 \\
  \texttt{email@domain} \\}

%\author{
%  \textbf{First Author\textsuperscript{1}},
%  \textbf{Second Author\textsuperscript{1,2}},
%  \textbf{Third T. Author\textsuperscript{1}},
%  \textbf{Fourth Author\textsuperscript{1}},
%\\
%  \textbf{Fifth Author\textsuperscript{1,2}},
%  \textbf{Sixth Author\textsuperscript{1}},
%  \textbf{Seventh Author\textsuperscript{1}},
%  \textbf{Eighth Author \textsuperscript{1,2,3,4}},
%\\
%  \textbf{Ninth Author\textsuperscript{1}},
%  \textbf{Tenth Author\textsuperscript{1}},
%  \textbf{Eleventh E. Author\textsuperscript{1,2,3,4,5}},
%  \textbf{Twelfth Author\textsuperscript{1}},
%\\
%  \textbf{Thirteenth Author\textsuperscript{3}},
%  \textbf{Fourteenth F. Author\textsuperscript{2,4}},
%  \textbf{Fifteenth Author\textsuperscript{1}},
%  \textbf{Sixteenth Author\textsuperscript{1}},
%\\
%  \textbf{Seventeenth S. Author\textsuperscript{4,5}},
%  \textbf{Eighteenth Author\textsuperscript{3,4}},
%  \textbf{Nineteenth N. Author\textsuperscript{2,5}},
%  \textbf{Twentieth Author\textsuperscript{1}}
%\\
%\\
%  \textsuperscript{1}Affiliation 1,
%  \textsuperscript{2}Affiliation 2,
%  \textsuperscript{3}Affiliation 3,
%  \textsuperscript{4}Affiliation 4,
%  \textsuperscript{5}Affiliation 5
%\\
%  \small{
%    \textbf{Correspondence:} \href{mailto:email@domain}{email@domain}
%  }
%}

\begin{document}
\maketitle
\begin{abstract}
\end{abstract}

\section{Introduction}

\begin{itemize}
    \item Documents are the primary medium through which organisations record and communicate information, with many critical real-world document types such as contracts, financial reports, scientific papers, technical manuals, textbooks, inherently spanning multiple pages and containing interleaved text, tables, figures and complex layouts \cite{tito2023hierarchical, borchmann2025arctic}
    \item Manually extracting information is time-consuming, expensive, and prone to human error; development of automated DocVQA systems that can answer natural-language queries over document images aims to solve this
    \item Significant progress has been made on single-page docVQA, with current SOTA models on-par if not exceeding human level performance, but vast majority of real-world documents extend beyond single page \cite{tito2023hierarchical}
    \item Evidence required to answer a question is often located on multiple pages scattered throughout the document, requiring a model to navigate, retrieve and reason across the entire document \cite{blau2024gram, borchmann2025arctic}
    \item challenges from single-page to multi-page VQA: 1. input token count scales linearly (or worse) quickly exceeding context window of existing LLMs and MLLMs, visually rich pages containing high resolution figures/tables/charts contribute to this. 2. capturing cross-page semantic and logical dependencies requires explicit long-range reasoning. 3. Explainability becomes more difficult as models need to also pinpoint how they arrived at the answer via answer page/passage retrieval.
    \item Early mult-page approaches adapt pre-trained language models to document domain treating OCR'd documents as a long text sequence however these methods ignored visual content.
    \item Subsequent work introduced hierarchical page-level architectures \cite{tito2023hierarchical} and global-reasoning token mechanisms \cite{blau2024gram} to enable cross-page reasoning and interaction achieving strong results but struggled with longer documents
    \item Unlike existing surveys that focus on single-page and broad document VRDU \cite{ding2025survey}, this survey focuses exclusively on the multi-page frontier for document question answering.
    \item we survey 30 (could be more) representative models spanning both OCR-dependent and OCR-free approaches.
\end{itemize}

\section{Framework Architecture Overview}
Multi-page VRDU frameworks can be compared along two related axes:
\begin{enumerate}
    \item OCR Usage: some frameworks rely on on external or localised OCR or PDF parsing as a primary input, while others read directly from document images using vision-language backbone.
    \item Document-level Architecture: all multi-page frameworks apply document-level architecture to control the growth in pages, tokens and evidence locations in long documents.
\end{enumerate}
Table~\ref{tab:model_survey_arch} summarises the surveyed frameworks along these axes.

% ─────────────────────────────────────────────────────────────────────────────
% Auto-generated by summary_tables_convert.py — do NOT edit manually.
% Re-generate:  python3 summary_tables_convert.py models_summary.csv models_summary.tex
%
% Required packages (add to your preamble):
%   \usepackage{booktabs}
%   \usepackage{longtable}
%   \usepackage{array}
%   \usepackage{xcolor}
%   \usepackage{colortbl}
%   \usepackage{caption}
%   \usepackage{adjustbox}   % for \begin{adjustbox}
%   \usepackage{natbib}      % for \citeyear  (or use biblatex)
%   \usepackage{lscape}      % optional: landscape
% ─────────────────────────────────────────────────────────────────────────────

\definecolor{sectionbg}{gray}{0.88}

% ── tab:model_survey_arch ──────────────────────────────────────────────────────────────
\begin{table*}[t]
\centering
% \small
\scriptsize
\captionsetup{font=small}
\setlength{\tabcolsep}{3pt}
\renewcommand{\arraystretch}{0.86}
\begin{adjustbox}{width=\textwidth}
\begin{tabular}{l l c p{2.8cm} p{3.0cm} c c c c}
\toprule
\textbf{Model} & \textbf{Venue} & \textbf{OCR} & \textbf{Vision Encoder} & \textbf{LLM Backbone} & \textbf{Mod.} & \textbf{Prompt} & \textbf{RAG} & \textbf{Agentic} \\
\midrule

% ── Hierarchical Document Transformer ─────────────────────────────────────────────────────────────
\rowcolor{sectionbg}
\multicolumn{9}{l}{\textbf{Hierarchical Document Transformer}} \\
\midrule
  \textbf{Arctic-TILT}~(\citeyear{borchmann2025arctic}) & ACL & Yes & U-Net & T5-Large & T, V, L & - & - & - \\
  \textbf{GRAM}~(\citeyear{blau2024gram}) & CVPR & Yes & DocFormerv2 & DocFormerv2 & T, V, L & - & - & - \\
  \textbf{Hi-VT5}~(\citeyear{tito2023hierarchical}) & PR & Yes & DiT & T5-Base & T, V, L & - & - & - \\
  \textbf{RM-T5}~(\citeyear{dong2024multi}) & DAS & Yes & DiT & T5-Base & T, V, L & - & - & - \\

\midrule[0.4pt]

% ── Backbone-Centric MLLM Adaptation ─────────────────────────────────────────────────────────────
\rowcolor{sectionbg}
\multicolumn{9}{l}{\textbf{Backbone-Centric MLLM Adaptation}} \\
\midrule
  \textbf{DocR1}~(\citeyear{xiong2026docr1}) & AAAI & No & Qwen2.5-VL & Qwen2.5-VL-7B & V & CoT & - & - \\
  \textbf{Leopard}~(\citeyear{jia2024leopard}) & TMLR? & No & SigLIP-SO400M & LLaMa-3.1-8B, Mistral-7B & V & CoT & - & - \\
  \textbf{mPLUG-DocOwl2}~(\citeyear{hu2025mplug}) & ACL & No & ViT-L/14 & mPLUG-Owl2 & V & - & - & - \\
  \textbf{Docopilot}~(\citeyear{duan2025docopilot}) & CVPR & No* & InternViT-300M & Multiple (Intern) & V & - & - & - \\
  \textbf{Texthawk2}~(\citeyear{yu2024texthawk2}) & preprint & No* & SigLIP-SO400M & Qwen2-7B-Instruct & V & - & - & - \\
  \textbf{DocSLM}~(\citeyear{hannan2025docslm}) & CVPR? & Yes & SigLIP2 & Qwen2.5-1.5B & T, V, L & - & - & - \\
  \textbf{InstructDr}~(\citeyear{tanaka2024instructdoc}) & AAAI & Yes & CLIP & FlanT5, BLIP2 & T, V, L & - & - & - \\
  \textbf{LayTokenLLM}~(\citeyear{zhu2025simple}) & CVPR & Yes & - & LLaMA3-8B, Qwen1.5-7B & T, L & - & - & - \\
  \textbf{CoR}~(\citeyear{guo2026end}) & preprint & Yes* & Qwen2.5-VL & Qwen2.5-VL-7B & V & CoT & - & - \\

\midrule[0.4pt]

% ── Retriever-Generator ─────────────────────────────────────────────────────────────
\rowcolor{sectionbg}
\multicolumn{9}{l}{\textbf{Retriever-Generator}} \\
\midrule
  \textbf{AVIR}~(\citeyear{li2025avir}) & ACM & No & Pix2Struct & Qwen2.5-VL-3B-AWQ & V & - & Dense & - \\
  \textbf{DFVC}~(\citeyear{qian2026accurate}) & MDPI & No & VisRAG, ColQwen & Qwen2-VL-2B-Instruct & V & - & Dense & - \\
  \textbf{M3DocRAG}~(\citeyear{cho2024m3docrag}) & preprint & No & ColPali, ColQwen & Multiple (Qwen/Intern/Idefics) & V & - & Dense & - \\
  \textbf{MoLoRAG}~(\citeyear{wu2025molorag}) & EMNLP & No & ColPali & Multiple (Qwen/DS/LLaVA) & V & - & Dense & - \\
  \textbf{Self-Attention Scoring}~(\citeyear{kang2024multi}) & ICDAR & No & Pix2Struct & Pix2Struct & V & - & Dense & - \\
  \textbf{VDocRAG}~(\citeyear{tanaka2025vdocrag}) & CVPR & No* & Phi3V-4.2B & Phi3V-4.2B & V & - & Dense & - \\
  \textbf{CREAM}~(\citeyear{zhang2024cream}) & ACM & Yes & Pix2Struct & LLaMA2-7B & T, V & - & Joint & - \\
  \textbf{MHier-RAG}~(\citeyear{gong2025mhier}) & preprint & Yes & Qwen-VL-Plus & Multiple (Qwen/DS/GPT) & T, V, L & CoT & Joint & - \\
  \textbf{MLDocRAG}~(\citeyear{zhang2026mldocrag}) & preprint & Yes & - & Multiple (Qwen) & T, V, L & CoT & Dense & - \\
  \textbf{MultiDocFusion}~(\citeyear{shin2025multidocfusion}) & ACL & Yes & - & Multiple (Qwen/LLaMa/Mistral) & T, V, L & - & Dense & - \\
  \textbf{PDF-WuKong}~(\citeyear{xie2024wukong}) & preprint & Yes* & IXC2-VL-4KHD & IXC2-VL-4KHD & T, V & - & Joint & - \\
  \textbf{RAG-DocVQA}~(\citeyear{lopez2025enhancing}) & ICDAR & Yes* & DiT, Pix2Struct & Qwen2.5-VL-7B-Instruct & T, V, L & - & Dense & - \\

\midrule[0.4pt]

% ── Agentic Pipeline ─────────────────────────────────────────────────────────────
\rowcolor{sectionbg}
\multicolumn{9}{l}{\textbf{Agentic Pipeline}} \\
\midrule
  \textbf{Doc-V*}~(\citeyear{zheng2026doc}) & preprint & No & ColQwen & Qwen2.5-VL-7B-Instruct & V & ReAct & Dense & Open \\
  \textbf{DocReact}~(\citeyear{wu2025doc}) & ACL & No & ColPali, VisRAG & GPT-4o & V & ReAct & Dense & Open \\
  \textbf{KGP}~(\citeyear{wang2024knowledge}) & AAAI & No & - & Multiple & T, L & - & Sparse & Fixed \\
  \textbf{MACT}~(\citeyear{yu2025visual}) & preprint & No & Multiple & Multiple (Qwen/Intern) & T, V & CoT & Dense & Open \\
  \textbf{DocLens}~(\citeyear{zhu2025doclens}) & preprint & Yes & - & Multiple (Closed-Source) & T, V, L & CoT & Joint & Open \\
  \textbf{DocAgent}~(\citeyear{sun2025docagent}) & EMNLP & Yes* & - & Multiple (Closed-Source) & T, V, L & ReAct & Joint & Open \\
  \textbf{DocDancer}~(\citeyear{zhang2026docdancer}) & preprint & Yes* & Qwen3-VL-235B & Qwen3-4B/30B-A3B & T, V, L & ReAct & Joint & Open \\
  \textbf{M2RAG}~(\citeyear{duan2025m2rag}) & NeurIPS & Yes* & VisRAG-Ret & Multiple (Qwen) & T, V & - & Joint & Fixed \\
  \textbf{MDocAgent}~(\citeyear{han2025mdocagent}) & preprint & Yes* & ColPali, ColQwen2 & Multiple (Llama/Qwen) & T, V & - & Joint & Fixed \\
  \textbf{SimpleDoc}~(\citeyear{jain2025simpledoc}) & EMNLP & Yes* & ColQwen-2.5 & Multiple (Qwen) & V & CoT & Joint & Open \\
  \textbf{ViDoRAG}~(\citeyear{wang2025vidorag}) & EMNLP & Yes* & NV-Embed-V2, ColQwen2 & Multiple (Qwen/LLaMa/GPT) & T, V & ReAct & Joint & Open \\

\bottomrule
\end{tabular}
\end{adjustbox}
\caption{Survey of VRDU models grouped by architecture category. \textbf{Mod.}: input modality (T\,=\,text, V\,=\,visual, L\,=\,layout). \textbf{OCR}: whether the model requires OCR pre-processing (Yes\,=\,required; No\,=\,not required; Yes\textsuperscript{*}\,=\,not strictly OCR-dependent but OCR is used as a framework component; No\textsuperscript{*}\,=\,OCR used during training only, not at inference). \textbf{RAG}: retrieval strategy (Dense, Sparse, Joint\,=\,both). \textbf{Prompt}: prompting strategy (CoT\,=\,Chain of Thought; ReAct\,=\,Reason\,+\,Act). \textbf{Agentic}: agentic workflow type (Open\,=\,open-ended tool use; Fixed\,=\,fixed pipeline). `--' indicates not reported.}
\label{tab:model_survey_arch}
\end{table*}


% Note sure whether to write "framework", "model" or "System"
\subsection{OCR Dependency}
The distinction between OCR-dependent and OCR-free document models was clearer in earlier VRDU works, where models either consumed OCR tokens and layout boxes or attempted end-to-end image-based reading \cite{huang2022layoutlmv3, ding2025survey}. In MP-VRDU, this binary taxonomy is less stable. Many models use OCR as one component inside a larger multimodal architecture rather than as the sole source of document evidence. OCR may provide sparse lexical retrival, page grounding, bounding-box references, or text snippets for verification, while images remain available for visual reasoning \cite{han2025mdocagent, zhu2025simple, wang2025vidorag}.

\paragraph{OCR Frameworks.}
OCR remains useful as structured text and layout coordinates give a compact signal for text-heavy or low-resolution pages. Layout-aware models such as LayTokenLLM \cite{zhu2025simple} show that OCR tokens, bounding boxes, and layout cues can be projected into a language-model space for efficient document reasoning, reducing the need to encode every page at high visual resolution. This however exposes the pipeline to OCR error propagation and parser failures \cite{liao2025doclayllm}. In multi-page settings, OCR-based models require retrieval, compression, or aggregation mechanisms \cite{hannan2025docslm, zhang2024cream}.

% \begin{itemize}
%     \item OCR-dependent frameworks extract text and often layout information via bounding-box coordinates using an OCR Engine or PDF parser. These tools can either be off the shelf or localised (self train transformer to perform OCR, often using SFT OCR objective).
%     \item The principle advantage is that clean, structured text provides a strong signal for reasoning at low image resolution, effectively reducing the reliance on high resolution document images 
%     \item Some frameworks \cite{zhu2025simple, wang2024docllm} perform inference purely on OCR-extracted information, yielding strong performance on text-heavy documents but lacking reasoning capabilities over complex visual elements.
%     \item As a result, a common architecture feeds OCR output, including the OCR text, 2D bounding-box coordinates, and low-resolution images patches, into a pre-trained document encoder \cite{huang2022layoutlmv3, appalaraju2024docformerv2}.
%     \item The fused representation is projected into an LLM's input space via an MLP or cross-attention adaptor. \cite{zhu2025simple}
%     \item This enables clean structured text to be extracted even from low-resolution scans, reducing the visual resolution required by the vision encoder, but is fundamentally limited by OCR error propagation \cite{wang2024docllm, liao2025doclayllm}.
%     \item In the multi-page setting, the primary challenge is that the OCR token count grows linearly with the number of pages, rapidly saturating the LLM context window.
%     \item multi-page OCR-dependent frameworks therefore require additional mechanisms for token reduction or cross-page information aggregation.
%     \item Due to the complexity of multi-page VRDU, OCR is often just one component of the overall architecture, and many models \cite{han2025mdocagent, jain2025simpledoc, wang2025vidorag} are not dependent on OCR, but use OCR for grounding and/or improving retrieval/agentic performance.
% \end{itemize}

\paragraph{OCR-Free Frameworks.}
OCR-free models encode document images directly and learn text recognition through visual-language pre-training \cite{lee2023pix2struct, hu2024mplug}. This preserves typography, spatial arrangement, figures, tables, and other visual details that OCR may discard. However, accurate image-based reading requires high-resolution document images, the cost of which grows rapidly as the number of pages scale upwards. Multi-page OCR-free frameworks therefore combine dynamic resolution, visual token compression, sparse sampling or iterative navigation to keep multi-page inference tractable \cite{yu2024texthawk2, zheng2026doc}.

% \begin{itemize}
%     \item OCR-free frameworks bypass text extraction by directly consuming high-resolution document images through a vision encoder.
%     \item Text recognition is instead learned through large-scale pre-training on document image–text pairs \cite{kim2021donut, radford2021learning, lee2023pix2struct}.
%     \item This enables end-to-end differentiability and more naturally preserves visual layout properties (font style, colour, spatial arrangement) that OCR struggles with or discards.
%     \item However, for accurate text recognition to be possible, a large number of visual tokens are produced for high-resolution input, which is manageable in single-page settings, becomes increasingly expensive when scaling to the multi-page setting. \cite{hu2025mplug}
%     \item effective text recognition also requires substantial pre-training on document specific tasks such as image captioning and masked image modelling, imposing significant datase and compute requirements. \cite{radford2021learning, li2022dit}
%     \item Representative single-page models include mPLUG-docOWL1.5 \cite{hu2024mplug} and Qwen2.5-VL \cite{bai2025qwen2} which combine dynamic resolution cropping with token compression modules to reduce per-image token conuts
%     \item In the multi-page setting, visual token count increases as document length increases as visual tokens are typically far more numerous than OCR text tokens for the same document length, making the context window problem significantly harder in OCR-free pipelines \cite{hannan2025docslm}
%     \item OCR-free multi-page frameworks need dedicated strategies for compression, retrieval or iterative document navigation to be feasible.
% \end{itemize}

% Not sure if need to cite here.
\subsection{Multi-Page VRDU Architectures and Trends}
Architectural progress in MP-VRDU has broadly evolved from direct extension of single-page document models towards increasingly selective and iterative document-level systems. Early work focused on adapting single-page transformer based document encoders to multi-page inputs. Later approaches moved towards backbone-centric MLLMs trained or tuned for multi-page reasoning. As document length remained a persistent constraint, retrieval-augmented methods became more prominent. More recent models increasingly adopt agentic workflows, where document understanding is framed as an iterative process of evidence seeking, reasoning, and refinement.
% \begin{itemize}
%     % Lewei Note: info here somewhat overlaps with introduction, could use improvement
%     \item In the multi-page setting, frameworks must additionally address two interrelated problems:
%     \item 1. Token count can scale linearly (for frameworks that use vision encoder) as document length increases, rapidly saturating the LLM context window; performance also degrades well before the context window is saturated, as a result of attention dilution \cite{liu2024lost}.
%     \item 2. Evidence required to answer a question is often scattered across pages, requiring explicit cross-page information flow. Another consideration is explainability, not only do models need to find evidence scattered cross multiple pages, they also need to explicitly locate the evidence.
% \end{itemize}

\paragraph{Hierarchical and Document-Structured Transformer.}
This family of architectures represents an early trend in MP-VRDU: imposing an explicit page-to-document reasoning structure on top of single-page document models. They usually begin from a single-page document transformer or encoder-decoder and scale to the multi-page setting by encoding pages separately, then aggregating information through page summary tokens \cite{tito2023hierarchical}, document global tokens and local-global attention \cite{blau2024gram}, sparse attention \cite{borchmann2025arctic}, or recurrent memory \cite{dong2024multi}. The main strength of this family is that it preserves page structure while allowing document-level interaction inside the model, making it more principled than concatenating all pages into one flat sequence. However, these models process all or most pages jointly, so memory and computation continue to increase with document length. They improve corss-page reasoning, but very long documents still require additional sparsity, compression, or retrieval.
% note: could expand more on model specifics but keep short for now.

% \begin{itemize}
%     \item This family of approaches introduces an explicit page-to-document reasoning structure, typically by encoding each page separately and then aggregating information through page summary tokens, document-global tokens, local-global attention, or recurrent memory.
%     \item The appeal is that strong single-page document backbones can be extended to the multi-page setting in a principled way, while preserving page structure and enabling cross-page interactions inside the model itself.
%     \item Compared with naive page concatenation, these architectures are better suited to cross-page reasoning, but they still process all or most pages jointly, so computation and memory usage continue to grow with document length.
%     \item In practice, they improve multi-page reasoning without fully removing the long-context bottleneck; very long documents remain difficult unless additional sparsity or compression is introduced.
%     \item Representative models include Hi-VT5 \cite{tito2023hierarchical}, GRAM \cite{blau2024gram}, RM-T5, and Arctic-TILT.
% \end{itemize}

\paragraph{Backbone-Centric MLLM Adaptation.}
Backbone-centric models represent a second trend: training or adapting a pre-trained LLM or MLLM to handle multi-page documents more directly, without adding a separate retrieval pipeline. They use mechanisms such as projectors, adapters, layout tokens, long-context tuning, reinforcement learning, streaming inference, or visual token compression to handle multi-page documents. This design allows the reuse of instruction following and multimodal reasoning ability of strong foundation models while making only target changes to the document understanding task. Some models \cite{jia2024leopard, tanaka2024instructdoc} rely on multi-page instruction tuning and the backbone context window to absorb page representations. Others \cite{hu2024mplug, hannan2025docslm, yu2024texthawk2} make compression or long-document reading central to the architecture. Recent work further adapts multimodal backbones toward document-level reasoning and evidence-guided learning \cite{duan2025docopilot, xiong2026docr1}. The limitation of this trend is that the backbone must still retain document-scale evidence internally, making performance sensitive to context length, compression quality and fine-grained text preservation.
 
% \begin{itemize}
%     \item This family adapts a pretrained LLM or MLLM backbone to multi-page documents through projectors, adapters, layout tokens, long-context tuning, reinforcement learning, streaming mechanisms, or visual compression modules, without introducing a separate document-level retrieval pipeline.
%     \item The appeal is that a capable foundation model can be reused directly, often with relatively modest architectural changes, while inheriting strong multimodal or instruction-following ability.
%     \item Some models rely mainly on multi-image instruction tuning and the backbone context window to absorb concatenated page representations, while others introduce native compression modules so that more pages can remain in context simultaneously.
%     \item The core limitation is that the backbone must still absorb document-scale evidence internally, so performance depends heavily on context length, compression quality, and the model's ability to retain fine-grained textual and layout detail.
%     \item Representative models include Leopard \cite{jia2024leopard}, InstructDoc \cite{tanaka2024instructdoc}, mPLUG-DocOwl2 \cite{hu2025mplug}, DocSLM \cite{hannan2025docslm}, Docopilot, and TextHawk2.
% \end{itemize}

\paragraph{Retriever-Generator.}
Retriever-Generator (RAG) models represent a third trend: using retrieval to reduce the amount of document evidence passed to the generator, decomposing MP-VRDU into evidence selection followed by answer generation. These systems first retrieve relevant pages, patches, chunks, visual elements, OCR passages or graph nodes, then generate an answer from the selected evidence. Retrieval may be performed through course-to-fine selection \cite{zhang2024cream}, adaptive visual retrieval \cite{li2025avir}, multimodal dense retrieval \cite{cho2024m3docrag, tanaka2025vdocrag}, sparse sampling \cite{xie2024wukong}, graph-based prompting \cite{wang2024knowledge}, or logic-aware retrieval \cite{wu2025molorag}. The main limitation of this architecture is the retrieval bottleneck. If the correct evidence is missed, the generator often cannot recover, regardless of its reasoning capability. 

% \begin{itemize}
%     \item Retrieve-then-generate frameworks decompose multi-page VRDU into an evidence selection stage followed by answer generation conditioned only on the selected pages, patches, chunks, diagrams, or graph nodes.
%     \item This sidesteps the token-scaling problem by making inference depend on the amount of retrieved evidence rather than the full document length.
%     \item The main drawback is the retrieval bottleneck: if the correct evidence page or passage is not retrieved, the generator will usually fail regardless of how strong the downstream MLLM or LLM is \cite{cho2024m3docrag, tanaka2025vdocrag, li2025avir}.
%     \item These systems can retrieve over OCR text, page images, multimodal chunks, or structured indexes such as graphs, and some models blur the boundary with native MLLMs by training the retriever or sampler end-to-end with the generator.
%     \item Representative models include CREAM \cite{zhang2024cream}, KGP \cite{wang2024knowledge}, AVIR \cite{li2025avir}, M3DocRAG \cite{cho2024m3docrag}, VDocRAG \cite{tanaka2025vdocrag}, and PDF-WuKong.
% \end{itemize}

\paragraph{Agentic Pipeline.}
Agentic pipelines represent the most recent trend: replacing fixed retrieve-then-generate workflows with iterative evidence seeking, reasoning and refinement. In this survey, we use this category for models with multi-agent decomposition and/or open-ended iterative control, rather than fixed RAG pipelines. These systems may distribute retrieval, visual inspection, modal fusion, judgement and answer synthesis across specialised agents \cite{han2025mdocagent, duan2025m2rag, yu2025visual}, or use tool-augment and dynamic iterative workflows to navigate documents, inspect evidence, maintain intermediate states, and revise answers \cite{sun2025docagent, zhu2025doclens, zheng2026doc, wang2025vidorag}. These models reduce single-shot retrieval failures and improve question answering accuracy, but introduce higher inference cost, latency and pipeline complexity.

% \begin{itemize}
%     \item Agentic frameworks extend retrieval-based reasoning by making evidence acquisition iterative and decision-driven rather than fixed in a single retrieve-then-generate pass.
%     \item A controller LLM or VLM decides what to inspect next based on intermediate observations, enabling query reformulation, page navigation, crop selection, verification, reflection, and multi-hop evidence gathering.
%     \item This can mitigate the single-shot retrieval bottleneck and is especially useful when evidence is distributed across multiple pages or modalities, but it increases inference cost, latency, and pipeline complexity.
%     \item Multi-agent variants further distribute work across specialized roles, such as text retrieval, visual retrieval, critical evidence extraction, answer synthesis, and adjudication, reducing the burden on any single model.
%     \item Representative models include MDocAgent \cite{han2025mdocagent}, DocAgent, DocLens, Doc-V*, M2RAG, ViDoRAG, and MACT.
% \end{itemize}

\section{Multimodal Information Modelling}

\subsection{Text Modality}
Unlike in single page documents, text passages can span multiple consecutive pages, relevant text can also be scattered across a number of pages in the document.

\paragraph{Auxillary Encoder as Structured Page-Level Representations}
\begin{itemize}
    \item (summarise corresponding section in general MLLM VRDU survey, use of encoders such as LayoutLMv3) 
    \item In multi-page setting, many models adapt single page auxillary encoder method for page-level representations. 
    \item Some models (GRAM, AVIR) use cross-attention on page representations for global reasoning. 
    \item Other models (hi-vt5, self-attention scoring, DocSLM) compress page representations then concatenate them for inference.
    \item Arctic-TILT uses sparse attention instead of concatentation to reduce quadratic attention to linear.
\end{itemize}

% Lewei Note: problem with this one is that it is not specific to text, also for visual modality, could exlude or keep it short.
\paragraph{Extracted Text as Long-Context / Retrieved Evidence}
\begin{itemize}
    \item OCR'd text is extracted and treated as a long textual corpus. 
    \item In multi-page setting, key problem is that the full document (even just the text) is too long and porrly ordered.
    \item Naive text chunks break section hiearchy, captions, tables etc.
    \item Perform page-aware chunking, heading trees, graph retrieval, parent-page retrieval, summary nodes, reranking, all done via LLM backbone.
    \item E.g. RAG-DocVQA, MDocAgent, M2RAG, MultiDocFusion
\end{itemize}

\paragraph{Text as One Modality in Hybrid Multimodal Retrieval and Agentic Pipelines}
\begin{itemize}
    \item RAG/Agentic models retrieve or reasons over textual and visual features. Text is useful as it is easy to search over, but shouldn't be fully trusted.
    \item Text retrieval may find keyword matching or semantically matching evidence but visually irrelevant pages.
    \item Many models separate text and image retrievers, then use inspector agents, modal fusers or cross-modal answer verification for retrieval or reasoning.
    \item E.g. CREAM, MDocAgent, M2RAG, DocLens (mainly RAG/agentic frameworks).
\end{itemize}

\paragraph{Text as a visually latent skill or training objective}
\begin{itemize}
    \item (summarise corresponding section in general MLLM survey). Text as training objective is still relevant in multi-page setting as text now spans multiple pages
    \item text can be learned from images through pretraining using OCR'd text, training data is synthetic (via OCR) so requires little human work.
    \item Training tasks include text reconstruction, grounding, layout-aware parsing.
    \item Problem in multi-page setting is requirement of high-resolution pages for this method, resulting in large number of visual tokens.
    \item Models use dynamic high-resolution tiling, vision token compression, page/section retrieval, course-to-fine visual fetching.
\end{itemize}

\subsection{Visual Modality}
Unlike single-page documents, multi-page frameworks need to manage a large number of visual tokens in addition to cross-page evidence.

\paragraph{Low resolution OCR-grounded visual augmentation}
\begin{itemize}
    \item Visual features are often used as course page, crop or patch evidence while OCR text and layout tokens carry fine grained information for answer generation.
    \item Effectively reduces visual token burden in multi-page documents, but susceptible to OCR/parser error propagation.
    \item E.g. Arctic-TILT, GRAM, RAG-DocVQA
\end{itemize}

\paragraph{High-Resolution Encoding and Token Compression and Course-to-Fine Access}
\begin{itemize}
    \item High resolution page images are encoded so that text, layout, tables, figures etc. can be read directly from pixels
    \item To deal with high token count, models use tiling, dynamic crops, resamplers, compressors, sparse samplers, thumbnail-to-crop visual actions
    \item E.g. mPLUG-DocOwl2, Docopilot, DocR1, Leopard
\end{itemize}

\subsection{Layout Modality}
%Word Bounding (textline bounding box) -> Reading Order-> Page ID / Correlation between caption/figure, subsection/section
In addition to layout within a page, need to account for document level layout e.g. Sections that span multiple pages, caption and figure may be on different pages, subsection from later pages of the document is linked to main section from earlier in the document.

% Lewei Note: already covered in general MLLM VRDU survey but still relevant, could keep very short
\paragraph{Token-Level Spatial Encoding from OCR/Layout Boxes}
\begin{itemize}
    \item OCR tokens paired with 2D coordinates, spatial embeddings or relative position biases, remains a basic way to encode layout, but naively concatenating all page-level layout tokens in multi-page setting causes sequence growth and weak cross-page structure.
    \item Models include Hi-VT5, RM-T5, IntructDr and Arctic-TILT
\end{itemize}

\paragraph{Layout as Document Structure}
\begin{itemize}
    \item Especially in retrieval-based multi-page systems, layout is handled more structurally than geometrically.
    \item Rather than using bounding boxes directly, document is parsed into layout aware units such as paragraphs, tables, figures, headings, etc.
    \item Popular formats include KGs, XML, Trees. 
    \item This method retains document level structure well, resulting in better retrieval performance and as a result better answer generation.
    \item Models include MultiDocFusion, MHierRAG, MLDocRAG
\end{itemize}

\paragraph{Implicit Layout Grounding}
\begin{itemize}
    \item OCR-free models handle layout implicitly through visual grounding, region localization, cropping heads or detection heads.
    \item Model learns where information is on the page and how to navigate it as a training objective, oftentime through reinforcement learning (don't remember which model)
    \item In multi-page setting, this often becomes course-to-fine navigation cross pages and regions e.g. CREAM
    \item Models include Doc-V* and mPLUG-DocOwl2
\end{itemize}

\subsection{Multimodal Fusion Overview}
Mutltimodal fusion involves using projectors/adapters to transform dimensionality of modalities to the same LLM backbone vector dimension. Fusion can be categorised into the stage in which fusion occurs:

\paragraph{Early (full concatenation) Fusion}
\begin{itemize}
    \item Text, layout and visual features are fused together before or during main encoder stage, usually at token or page level.
    \item Modalities are embedded together and interact via self-attention or cross-attention
    \item Good for precise alignment between modalities, but scales poorly as page count grows
    \item Representative models: Hi-VT5, RM-T5, InstructDr, GRAM
\end{itemize}

\paragraph{Late (Selective) Fusion}
\begin{itemize}
    \item Modalities are processed more independently first, and only combined after retrieval, ranking, summarisation or page selection. 
    \item SOme models retrieve OCR text chunks and page images separately, then fuse them only at evidence selection stage or final answer stage
    \item Scales well for multi-page setting as token overhead is effectively reduced by avoiding fusing whole document jointly.
    \item Model only fuses small set of candidate pages, crops, chunks or agent outputs
    \item Problem is modality alignment can be poor, and retrieval/agentic performance is harder when modalities are separate.
    \item models include MDocAgent, M2RAG, ViDoRAG, CREAM
\end{itemize}

\paragraph{Staged (hierarchical) Fusion}
\begin{itemize}
    \item Many frontier multi-age models implement hierarchical or staged fusion, where early fusion is used within a page and late/hierarchical fusion across pages. 
    \item Arguably most natural multi-page strategy, preservces strong local multimodal alignment while keeping cross-page reasoning scalable, but not many models use it.
    \item Early models Hi-VR5 and RM-T5 use this, later models such as Doc-V* use this as well.
\end{itemize}

\section{Training Strategies}

% Lewei Note: some overlap with MLLM survey, could exclude
\paragraph{Parameter efficienct tuning on frozen backbones}
Multi-page end-to-end training is very expensive, so many models take off-the-shelf vision encoders and LLM backbones and instead use loRA, projectors or retrieval heads to adapt to multi-page setting. They use strong pretrained LLM / LVLMs and specialise it to document retrieval or QA. 
LayTokenLLM, CREAM, MoLoRAG all do this.

\paragraph{Curriculum / Multi-Stage Training}
Either directly adapt single page model, then train on multi-page objectives but doesn't scale well.
Alternatively train with easier local tasks then move towards harder multi-page objectives (end-to-end training) for better performance and efficiency.
Representative models include DocSLM: image-OCR alignment -> image compression -> document compression -> instructuoin tuning, streaming abstention.
mPLUG-DocOWL2 also.

\paragraph{Self-Supervised Retrieval Training}
Main approach with multi-page is RAG / retrieval, many models train retrieval transformer using techniques such as contrastive learning e.g. VDocRAG uses self-supervised representation compression objectives. 

\paragraph{Reinforcement Learning For Reasoning}
Other main approach with multi-page is agentic iterative reasoning, many models optimize reasoning process e.g. DocR1 uses evidence guided GRPO, DocV* uses SFT + GRPO

\section{Training Free Strategies}
% Prompt Engineering, RAG and Agentic Framework

\paragraph{Chain of Thought / ReAct}
many models use chain or thought (self reflection) to iteratively find more relevant / better evidence for answer inference. 
Models include Doc-Agent, Doc-React, DocLens, all very similar is technique.

% Overlap with Layout modality section
\paragraph{Structure-aware retrieval + other retrieval methods}
Construct knowledge base (structured) from document, then run RAG or agentic iterative retrieval to improve performance. Use document structure explicitly at inference time, can be constructed using off the shelf PDF parser, allows for hierarchy aware inference. E.g. KGP, MoLoRAG (both graph based), MultiDocFusion.

Other models perform RAG by combining different retrieval cues such as text embeddings, visual embeddings, OCR text, summaries, graph neighbours etc. e.g. ViDoRAG, MDocAgent

% Overlap with agentic architecture in VRDU framework overview
\paragraph{Multi-Agent Decomposition}
Given a query, split taslk into multiple roles such as planner, retriever, judge, visual inspector etc. Role specialisation and intermediate checking increases performance without changing model weights. Agents are often just same LLM backbone but with different prompt template. Models include DocAgent, DocLens, MDocAgent, ViDoRAG.

\section{Datasets}

\section{Challenges}

\section*{Acknowledgments}

% Bibliography entries for the entire Anthology, followed by custom entries
%\bibliography{anthology,custom}
% Custom bibliography entries only
\bibliography{custom}

\appendix

\section{Example Appendix}
\label{sec:appendix}

% ─────────────────────────────────────────────────────────────────────────────
% Auto-generated by summary_tables_convert.py — do NOT edit manually.
% Re-generate:  python3 summary_tables_convert.py models_summary.csv models_summary.tex
%
% Required packages (add to your preamble):
%   \usepackage{booktabs}
%   \usepackage{longtable}
%   \usepackage{array}
%   \usepackage{xcolor}
%   \usepackage{colortbl}
%   \usepackage{caption}
%   \usepackage{adjustbox}   % for \begin{adjustbox}
%   \usepackage{natbib}      % for \citeyear  (or use biblatex)
%   \usepackage{lscape}      % optional: landscape
% ─────────────────────────────────────────────────────────────────────────────

\definecolor{sectionbg}{gray}{0.88}

% ── tab:model_survey_ocr ──────────────────────────────────────────────────────────────
\begin{table*}[t]
\centering
% \small
\scriptsize
\captionsetup{font=small}
\setlength{\tabcolsep}{3pt}
\renewcommand{\arraystretch}{0.86}
\begin{adjustbox}{width=\textwidth}
\begin{tabular}{l l l l l c c c c}
\toprule
\textbf{Model} & \textbf{Venue} & \textbf{Architecture} & \textbf{Vision Encoder} & \textbf{LLM Backbone} & \textbf{Mod.} & \textbf{Prompt} & \textbf{Search} & \textbf{Agentic} \\
\midrule

% ── No ─────────────────────────────────────────────────────────────
\rowcolor{sectionbg}
\multicolumn{9}{l}{\textbf{No}} \\
\midrule
  {DocR1}~(\citeyear{xiong2026docr1}) & AAAI & BC & Qwen2.5-VL & Qwen2.5-VL-7B & V & CoT & - & - \\
  {Leopard}~(\citeyear{jia2024leopard}) & TMLR? & BC & SigLIP-SO400M & LLaMa-3.1-8B, Mistral-7B & V & CoT & - & - \\
  {mPLUG-DocOwl2}~(\citeyear{hu2025mplug}) & ACL & BC & ViT-L/14 & mPLUG-Owl2 & V & - & - & - \\
  {Docopilot$^*$}~(\citeyear{duan2025docopilot}) & CVPR & BC & InternViT-300M & Multiple (Intern) & V & - & - & - \\
  {Texthawk2$^*$}~(\citeyear{yu2024texthawk2}) & preprint & BC & SigLIP-SO400M & Qwen2-7B-Instruct & V & - & - & - \\
  {AVIR}~(\citeyear{li2025avir}) & ACM & RAG & Pix2Struct & Qwen2.5-VL-3B-AWQ & V & - & Dense & - \\
  {DFVC}~(\citeyear{qian2026accurate}) & MDPI & RAG & VisRAG, ColQwen & Qwen2-VL-2B-Instruct & V & - & Dense & - \\
  {M3DocRAG}~(\citeyear{cho2024m3docrag}) & preprint & RAG & ColPali, ColQwen & Multiple (Qwen/Intern/Idefics) & V & - & Dense & - \\
  {MoLoRAG}~(\citeyear{wu2025molorag}) & EMNLP & RAG & ColPali & Multiple (Qwen/DS/LLaVA) & V & - & Dense & - \\
  {Self-Attention Scoring}~(\citeyear{kang2024multi}) & ICDAR & RAG & Pix2Struct & Pix2Struct & V & - & Dense & - \\
  {VDocRAG$^*$}~(\citeyear{tanaka2025vdocrag}) & CVPR & RAG & Phi3V-4.2B & Phi3V-4.2B & V & - & Dense & - \\
  {Doc-V*}~(\citeyear{zheng2026doc}) & preprint & AP & ColQwen & Qwen2.5-VL-7B-Instruct & V & ReAct & Dense & Open \\
  {DocReact}~(\citeyear{wu2025doc}) & ACL & AP & ColPali, VisRAG & GPT-4o & V & ReAct & Dense & Open \\
  {KGP}~(\citeyear{wang2024knowledge}) & AAAI & AP & - & Multiple & T, L & - & Sparse & Fixed \\
  {MACT}~(\citeyear{yu2025visual}) & preprint & AP & Multiple & Multiple (Qwen/Intern) & T, V & CoT & Dense & Open \\

\midrule[0.4pt]

% ── Yes ─────────────────────────────────────────────────────────────
\rowcolor{sectionbg}
\multicolumn{9}{l}{\textbf{Yes}} \\
\midrule
  {Arctic-TILT}~(\citeyear{borchmann2025arctic}) & ACL & HDT & U-Net & T5-Large & T, V, L & - & - & - \\
  {GRAM}~(\citeyear{blau2024gram}) & CVPR & HDT & DocFormerv2 & DocFormerv2 & T, V, L & - & - & - \\
  {Hi-VT5}~(\citeyear{tito2023hierarchical}) & PR & HDT & DiT & T5-Base & T, V, L & - & - & - \\
  {RM-T5}~(\citeyear{dong2024multi}) & DAS & HDT & DiT & T5-Base & T, V, L & - & - & - \\
  {DocSLM}~(\citeyear{hannan2025docslm}) & CVPR? & BC & SigLIP2 & Qwen2.5-1.5B & T, V, L & - & - & - \\
  {InstructDoc}~(\citeyear{tanaka2024instructdoc}) & AAAI & BC & CLIP & FlanT5, BLIP2 & T, V, L & - & - & - \\
  {LayTokenLLM}~(\citeyear{zhu2025simple}) & CVPR & BC & - & LLaMA3-8B, Qwen1.5-7B & T, L & - & - & - \\
  {CoR$^*$}~(\citeyear{guo2026end}) & preprint & BC & Qwen2.5-VL & Qwen2.5-VL-7B & V & CoT & - & - \\
  {CREAM}~(\citeyear{zhang2024cream}) & ACM & RAG & Pix2Struct & LLaMA2-7B & T, V & - & Joint & - \\
  {MHier-RAG}~(\citeyear{gong2025mhier}) & preprint & RAG & Qwen-VL-Plus & Multiple (Qwen/DS/GPT) & T, V, L & CoT & Joint & - \\
  {MLDocRAG}~(\citeyear{zhang2026mldocrag}) & preprint & RAG & - & Multiple (Qwen) & T, V, L & CoT & Dense & - \\
  {MultiDocFusion}~(\citeyear{shin2025multidocfusion}) & ACL & RAG & - & Multiple (Qwen/LLaMa/Mistral) & T, V, L & - & Dense & - \\
  {PDF-WuKong$^*$}~(\citeyear{xie2024wukong}) & preprint & RAG & IXC2-VL-4KHD & IXC2-VL-4KHD & T, V & - & Joint & - \\
  {RAG-DocVQA$^*$}~(\citeyear{lopez2025enhancing}) & ICDAR & RAG & DiT, Pix2Struct & Qwen2.5-VL-7B-Instruct & T, V, L & - & Dense & - \\
  {DocLens}~(\citeyear{zhu2025doclens}) & preprint & AP & - & Multiple (GPT/Gemini/Claude) & T, V, L & CoT & Joint & Open \\
  {DocAgent$^*$}~(\citeyear{sun2025docagent}) & EMNLP & AP & - & Multiple (GPT/Gemini/Claude) & T, V, L & ReAct & Joint & Open \\
  {DocDancer$^*$}~(\citeyear{zhang2026docdancer}) & preprint & AP & Qwen3-VL-235B & Qwen3-4B/30B-A3B & T, V, L & ReAct & Joint & Open \\
  {M2RAG$^*$}~(\citeyear{duan2025m2rag}) & NeurIPS & AP & VisRAG-Ret & Multiple (Qwen) & T, V & - & Joint & Fixed \\
  {MDocAgent$^*$}~(\citeyear{han2025mdocagent}) & preprint & AP & ColPali, ColQwen2 & Multiple (Llama/Qwen) & T, V & - & Joint & Fixed \\
  {SimpleDoc$^*$}~(\citeyear{jain2025simpledoc}) & EMNLP & AP & ColQwen-2.5 & Multiple (Qwen) & V & CoT & Joint & Open \\
  {ViDoRAG$^*$}~(\citeyear{wang2025vidorag}) & EMNLP & AP & NV-Embed-V2, ColQwen2 & Multiple (Qwen/LLaMa/GPT) & T, V & ReAct & Joint & Open \\

\bottomrule
\end{tabular}
\end{adjustbox}
\caption{Survey of VRDU models grouped by OCR dependency. Groups: \textbf{Yes}\,=\,OCR required at inference; \textbf{No}\,=\,OCR not required at inference. $^*$\,next to a model name has group-dependent meaning: \textbf{Yes\textsuperscript{*}}\,=\,model is not strictly OCR-dependent but incorporates OCR as a framework component; \textbf{No\textsuperscript{*}}\,=\,OCR was used explicitly during training but is not required at inference. \textbf{Architecture} abbreviations: HDT\,=\,Hierarchical Document Transformer; BC\,=\,Backbone-Centric MLLM Adaptation; Search\,=\,Retriever-Generator; AP\,=\,Agentic Pipeline. \textbf{Mod.}: input modality (T\,=\,text, V\,=\,visual, L\,=\,layout). \textbf{Search Strat.}: retrieval strategy (Dense, Sparse, Joint\,=\,both). \textbf{Prompt Strat.}: prompting strategy (CoT\,=\,Chain of Thought; ReAct\,=\,Reason\,+\,Act). \textbf{Agentic}: agentic workflow type (Open\,=\,open-ended tool use; Fixed\,=\,fixed pipeline). `--' indicates not reported.}
\label{tab:model_survey_ocr}
\end{table*}


\end{document}
